\section{Benchmark Trade Taxes}\label{app:tariffs}

Figures \ref{fig:taxcom1} to \ref{fig:taxpar2} summarize the commodity and partner structure of the taxes on imports into the four countries in the \texttt{g20} aggregation, in the style of \citet{LanzRutherford2016}.
They show import tariff rates of the four importing countries and export tax rates that the exporting regions charge on exports to the four countries, together with the tariff revenue of the four countries (in billion USD).
Averaged over all imports, the United States levy import tariffs of 1.1\% and face export taxes of their trading partners of 0.5\%.
For China, the corresponding figures are 1.9\% and 0.5\%, for South Africa 4.1\% and 1.4\%, and for Argentina 4.9\% and 1.3\%.
Protection is far from uniform.
In the United States, the highest rates apply to sugar, leather, dairy products, and wearing apparel; in China, to wheat, sugar, other cereals, and wool.
South Africa and Argentina levy rates of 15--32\% on wearing apparel, leather, and textiles.
Across partners, preferential agreements are visible: the United States levy almost no tariffs on imports from Mexico and Canada, Argentina none on imports from Brazil.
Tariff revenue amounts to 43.7 billion USD in the United States, 53.8 billion USD in China, 5.4 billion USD in South Africa, and 4.3 billion USD in Argentina.
It is concentrated on a few goods and partners.
Wearing apparel and leather account for 38\% of the tariff revenue of the United States, and motor vehicles for 14\% of the tariff revenue of China and 33\% of that of South Africa.
Imports from China account for 34--43\% of the tariff revenue of the United States, South Africa, and Argentina.

\begin{figure}[p!]
\centering
\input{jgea_new/figures/tariff_commodity_USA}
\vspace{0.1in}
\input{jgea_new/figures/tariff_commodity_CHN}
\pgfplotslegendfromname{tariffcom_legend}
\caption{Trade taxes on imports by commodity, United States and China (\GTAP 12, 2023; 23 most highly taxed commodities).}\label{fig:taxcom1}
\end{figure}

\begin{figure}[p!]
\centering
\input{jgea_new/figures/tariff_commodity_ZAF}
\vspace{0.1in}
\input{jgea_new/figures/tariff_commodity_ARG}
\pgfplotslegendfromname{tariffcom_legend}
\caption{Trade taxes on imports by commodity, South Africa and Argentina (\GTAP 12, 2023; 23 most highly taxed commodities).}\label{fig:taxcom2}
\end{figure}

\begin{figure}[p!]
\centering
\input{jgea_new/figures/tariff_txtm_USA}
\vspace{0.1in}
\input{jgea_new/figures/tariff_txtm_CHN}
\pgfplotslegendfromname{tariffpar_legend}
\caption{Trade taxes on imports by trading partner, United States and China (\GTAP 12, 2023).}\label{fig:taxpar1}
\end{figure}

\begin{figure}[p!]
\centering
\input{jgea_new/figures/tariff_txtm_ZAF}
\vspace{0.1in}
\input{jgea_new/figures/tariff_txtm_ARG}
\pgfplotslegendfromname{tariffpar_legend}
\caption{Trade taxes on imports by trading partner, South Africa and Argentina (\GTAP 12, 2023).}\label{fig:taxpar2}
\end{figure}
